
\chapter{Loyalty: Factions \& Religion}
Factions and Religion are two different mechanics that grant power and progression opportunities 
based on Loyalty to another character or society of characters.


\section{Loyalty}
By default, there are 3 levels of \textit{Loyalty}, 
from -1 (unfavorable) to +1 (favorable), with 0 being neutral.

\vspace{0.4em}

This may vary depending on the faction or religion, 
with more complex societies having additional loyalty tiers 
depending on the temperament and organization of the society.

\begin{Example}
  A society that is particularly hostile to outsiders, while having a low internal complexity or cohesion,
  may have more unfavorable tiers without increasing the number of favorable tiers (e.g., -5 to +1).

  \vspace{0.4em}

  A peaceful society that does not differentiate between various outsider statuses, 
  or are more forgiving of disloyalty, while also having a rich internal hierarchy of trust, 
  may have no unfavorable tiers and many favorable tiers (e.g., 0 to +5).
\end{Example}

Loyalty can be increased by performing favors and going on missions,
or by satisfying certain character conditions 
(e.g., based on perks \& abilities, or on immutable characteristics like Race/Species).

\vspace{0.4em}

Loyalty can be decreased by failure on missions or by gaining loyalty with adversarial factions.


\section{Factions}
\textit{Factions} are societies of characters organized around a common purpose.

\vspace{0.4em}

Locations, including Towns (see Chapter 9: Regions) and Mission locations (see Chapter 10: Adventures), 
may be controlled by one or more factions, 
who control access to the location and/or its amenities and effects.

\vspace{0.4em}

Factions may offer unique training (e.g., perks \& abilities), missions, or plot-related opportunities.


\newpage


\section{Religion}
\textit{Religion} is mechanically independent from but may overlap with factions;
meaning, factions may organize around a religion and provide benefits that synergize with or 
depend on the perks and abilities of those religions (i.e., as prerequisites).

\begin{DesignerNote}
  This rulebook provides a framework for storytelling in fictional polytheistic worlds,
  but does not specify the rules for individual religions within those worlds.

  \vspace{0.4em}

  As such, it is not required to implement religion mechanics within specific campaigns 
  if it would cause scruples for the players.
\end{DesignerNote}


\subsection{Deity}
The \textit{Deity} of a religion may grant unique Perks and Abilities to their followers,
or may provide temporary buffs from \textit{Prayer} when a follower levels-up in a Town with a 
Temple dedicated to that deity.

\vspace{0.4em}

The specific perks granted by a religion may be contingent on the character's loyalty to that deity,
being introduced at different loyalty tiers.


\subsection{Death}
The \textit{Death} mechanic of a religion describes what happens to a character who follows that religion
when that character's Life is reduced to 0.

\begin{Example}
  A death mechanic may preserve a character in an afterlife, 
  so that their knowledge and wisdom may be relied upon by living characters,
  or even so that they may be returned to play under certain conditions.

  \vspace{0.4em}

  A more moderate (and familiar) example may preserve characters at 0 Life
  by having them become Unconscious, then force them to succeed on a number of 
  ``Death Saves'' (\Roll{d10}) in order to maintain that state.
\end{Example}


